%--- iliad ---
% separateSolutions: true
% title: A hierarchical clustering example
% summary: >-
%   Thirty-six labelled observations, shared location and shape parameters,
%   and three ways to organize the same observable law. A finite example of
%   residual uncertainty and the conditioned score.
%--- end ---
\documentclass[11pt]{article}
\usepackage[T1]{fontenc}
\usepackage{lmodern}
\IfFileExists{iliad.sty}{\usepackage[nogeometry]{iliad}}{\usepackage[nogeometry]{../iliad}}
\usepackage[a4paper,margin=22mm]{geometry}
\usepackage{microtype}
\title{A hierarchical clustering example}
\author{Satya Benson}
\date{September 21, 2026}
\begin{document}
\maketitle

\section{Prerequisites}
This supplementary example uses the latent-model definition and conditional
entropy from the main worksheet. All variables are finite-valued and all logs
have base two. The example adapts the hierarchy in Jeremy Gillen and Daniel
Chiang's \emph{A summary of Condensation and its relation to Natural Latents},
\S1.1.1. The construction and figure here are new; we do not use their numerical
entropy estimates.

\begin{learningoutcomes}
\begin{itemize}
\item Distinguish LVM, where the latents assigned to an observation determine
it, from reconstruction, where an observation recovers those latents.
\item Compare the information charged by different latent models for the same query.
\end{itemize}
\end{learningoutcomes}

\section{Thirty-six observations}
There are two groups, each containing two clusters, each with nine observations.
Write $X_{gcj}$ for the $j$th two-dimensional point in cluster $c$ of group $g$.
The labels $(g,c,j)$ and memberships are fixed. We are comparing models of this
labelled family, not inferring an unknown partition of an unlabelled cloud.

\begin{figure}[ht]
\centering
\includegraphics[width=\linewidth]{fig/hierarchical-points.pdf}
\caption{One draw from the finite model. The left panel shows the 36 points;
the right panel shows which set of observations receives each parameter, and is
not a graph inferred from the points. The global shape $S$ is available to every
observation.}
\label{fig:hierarchy}
\end{figure}

The parameters are a global shape $S$, group centers $G_1,G_2$, and cluster
centers $M_{gc}$. Given these parameters, the 36 observations are independent.
Marginalizing over shared parameters generally makes observations dependent.
The plotted cloud is one realization; the random variables and their entropies
refer to the joint law across realizations.

\newpage
\section{An explicit finite law}
Let $D=\{0,\ldots,200\}^2$. Draw $S$ uniformly from three shape labels, and
draw $G_1,G_2$ independently and uniformly from $\{60,\ldots,140\}^2$,
independently of $S$. Conditional on $G_g$, draw its two centers independently:
\[
 P(M_{gc}=m\mid G_g=a)\ \propto\
 \exp\!\left(-\frac{\|m-a\|^2}{128}\right),\qquad m\in D.
\]
Given $S$ and all centers, draw the observations independently, using
\[
 P(X_{gcj}=x\mid M_{gc}=m,S=s)\ \propto\
 \exp\!\left(-\tfrac12(x-m)^\mathsf{T}Q_s^{-1}(x-m)\right),\quad x\in D,
\]
where
\[
 Q_{\mathrm{round}}=\begin{pmatrix}5&0\\0&5\end{pmatrix},\quad
 Q_{\mathrm{right}}=\begin{pmatrix}25&20\\20&25\end{pmatrix},\quad
 Q_{\mathrm{left}}=\begin{pmatrix}25&-20\\-20&25\end{pmatrix}.
\]
Each proportionality means normalization by the sum over $D$. These are
probability mass functions on a finite grid. The matrices control shape; they
are not the exact covariances after discretization and boundary truncation.
Every grid point has positive probability. Cluster centers tend to lie near
their group center, but the four components need not be visually separated in
every realization.

\section{A structured latent model}
Let $I$ contain all 36 indices. Let $A_g$ be the 18 indices of group $g$, and
$A_{gc}$ the nine indices of cluster $(g,c)$. Assign
\[
 Y_I=S,\qquad Y_{A_g}=G_g,\qquad Y_{A_{gc}}=M_{gc},\qquad
 Y_{\{i\}}=X_i,
\]
and make every other slot constant. The joint law factorizes in the address
inclusion order, so its ordered Markov defect is zero.

The singleton slot supplies the sampled point itself. Thus
$H\bigl(X_i\mid Y_{\{A\,:\,i\in A\}}\bigr)=0$ exactly. The cluster center and shape alone
specify a conditional distribution for the point; they do not determine its
realized coordinates. The residual conditional entropy is
$H(X_i\mid M_{gc},S)$.

Reconstruction is different. All the kernels have full support, so any
finite observation block leaves several parameter configurations with positive
posterior probability. For example, $H(M_{gc}\mid X_i)>0$. The model's joint
law satisfies the Markov condition exactly, but it is not a perfect condensation.

For a nonempty block $B$, let $\mathcal J_B=\{A:A\cap B\ne\varnothing\}$.
The general score-gap identity is
\[
 \chi(B)-H(X_B)
 =\bigl[\chi(B)-H(Y_{\mathcal J_B})\bigr]
   +H(Y_{\mathcal J_B}\mid X_B).
\]
Here the bracketed Markov term is zero. The gap is the latent information
remaining uncertain after observing the queried block.

\newpage
\section{Three organizations of the same observable law}
For a nonempty queried \emph{observation block} $B$, the conditioned score
sums the charges of all slots whose addresses meet $B$.

\paragraph{One global tuple.}
Put $Z_I=X_I$ and make every other slot constant. The global slot determines
each observation. For every nonempty $B$,
\[
 \chi_{\mathrm{global}}(B)=H(X_I),\qquad
 \chi_{\mathrm{global}}(B)-H(X_B)=H(X_I\mid X_B).
\]
This attains the entropy bound when querying the whole tuple. Even a singleton
query incurs the entire tuple's entropy in this score.

\paragraph{An ordered chain.}
Choose an ordering $X_1,\ldots,X_{36}$. Set
$Z_{\{j,j+1,\ldots,36\}}=X_j$ and make other slots constant. The contributing
slots for $X_i$ contain $X_1,\ldots,X_i$. The charge for slot $j$ is
$H(X_j\mid X_1,\ldots,X_{j-1})$, hence
\[
 \chi_{\mathrm{chain}}(B)=H(X_1,\ldots,X_{\max B}).
\]
This attains the entropy bound for prefix queries. Querying only $X_{36}$ costs
$H(X_I)$. These statements concern this particular address assignment, not
every autoregressive architecture.

\paragraph{The structured model.}
The score charges the shape once, the group and cluster parameters needed by
$B$, and the residual uncertainty of each requested observation:
\[
\begin{aligned}
 \chi_{\mathrm{structured}}(B)={}&H(S)
 +\sum_{g:B\cap A_g\ne\varnothing}H(G_g)\\
 &+\sum_{g,c:B\cap A_{gc}\ne\varnothing}H(M_{gc}\mid G_g)
 +\sum_{i\in B}H(X_i\mid M_{g(i)c(i)},S).
\end{aligned}
\]
This expression avoids charging for unrelated observations. It need not be
smaller than the other scores for every block. In particular, for $B=I$ it
exceeds $H(X_I)$ by the remaining parameter uncertainty, whereas both baseline
constructions attain $H(X_I)$. Fitting the observable law does not by itself
settle which organization is preferable across queries.

\section{Why blocks help recover parameters}
For nonempty $B$ inside a fixed cluster $(g,c)$, set $\Theta=(S,G_g,M_{gc})$. The
nonconstant contributing tuple is $(\Theta,X_B)$, so
\[
 \chi_{\mathrm{structured}}(B)-H(X_B)=H(\Theta\mid X_B).
\]
For nonempty $B\subseteq B'\subseteq A_{gc}$, the reduction in this gap is
\[
 H(\Theta\mid X_B)-H(\Theta\mid X_{B'})
 =I(\Theta;X_{B'\setminus B}\mid X_B)\ge0.
\]
More points can inform the shared parameters. Exact recovery still fails for
finite blocks. Even knowing one cluster center need not determine its parent
group center. When $B$ extends to another cluster, its contributing parameter
family changes, so the displayed monotonicity argument no longer applies as
written. This example motivates block witnesses and nonzero recovery errors;
it does not establish a chosen uniform error bound.

\paragraph{Source.}
The motivating hierarchy is from Gillen and Chiang,
\href{https://www.alignmentforum.org/posts/agw7HhW4cWjADpBgo}{\emph{A summary of
Condensation and its relation to Natural Latents}}, \S1.1.1. The finite law,
illustration, and exact score comparisons above are this course's adaptation.
\end{document}
